\section{Evaluation}
\label{sec:eval}
\label{sec:evaluation}

To empirically evaluate relational \ematching,
 we implemented it inside the
 \egg equality saturation toolkit~\cite{egg}.
Our implementation consists
 about 80 lines of Rust inside \egg itself to convert
 patterns into conjunctive queries,
 paired with a
 a separate, e-graph-agnostic
 Rust library to implement generic join in fewer than $500$ lines.

\egg's existing
 \ematching infrastructure is also about $500$ lines of Rust,
 and it is interconnected to various other parts of \egg.
Qualitatively,
 we claim that the relational approach is simpler to implement,
 especially since the CQ solver is completely modular.
 We could plug in a different generic join implementation~\footnote{
   There is no reusable generic join implementation at the time of writing.},
 or even a more conventional binary join implementation.

In this section, we refer to \egg's existing \ematching implementation as ``\EM''
 and our relation approach as ``\GJ.''

\subsection{Benchmarking Setup}

We use \egg's two largest benchmark suites
 as the basis for our two benchmark suites.
The \texttt{math} suite implements a
 simple computer algebra system,
 including limited support for symbolic differentiation and integration.
The \texttt{lambda} suite implements a
 partial evaluator for the lambda calculus.
Each \egg benchmark suite provides a set of rewrite rules,
 each with a left and right side pattern,
 and a set of starting terms.

To construct the \egraphs used in our benchmarks
 we ran equality saturation
 on a set of terms selected from \egg's test suite,
 stopping before the \egraph reached
 1e5, 1e6, 2e6, and 3e6 \enodes.
The result is four increasingly large \egraphs
 for each benchmark suite
 filled with terms that are
 generated by the suite's rewrite rules.
For each benchmark suite
 and each of the four \egraph sizes,
 we then ran \ematching on the \egraph
 using both \EM and \GJ.
We ran each approach 10 times and
 took the minimum run time.

For our \GJ approach,
 we ran each trial twice.
The first time builds the index tries
 necessary for generic join just-in-time,
 and the run time includes that.
On the second trial,
 \GJ uses the pre-built index tries from the first run,
 so the time to build them is excluded.
In both \autoref{fig:speedup} and \autoref{tab:eval},
 orange bars/rows show the first runs (including indexing),
 and blue bars/rows show the second runs (excluding indexing).

All benchmarks are single-threaded,
 and they were executed on a
 4.6GHz CPU with 32GB of memory.

\begin{figure}
  \centering
  \includegraphics[width=1.0\linewidth]{speedup.pdf}
  \caption{
    Relational e-matching can be up to 6 orders of magnitude faster
     than traditional e-matching on complex patterns.
    Speedup tends to be greater when the output size is smaller.
    Bars to the right of the ``$1\times$'' line
     indicate that relational \ematching is faster.
    The plots show two benchmark suites,
     \texttt{lambda} and \texttt{math},
     taken from the \egg test suite.
    Each group of bars
     shows the benchmarking results of e-matching a single
     pattern on 4 increasingly large e-graphs (top to bottom),
     comparing
     \egg's built-in e-matching (\EM)
     with our relational e-matching approach using generic join (\GJ).
    % Patterns are sorted
    %  by the number of \ematching results on the largest \egraph
    %  (shown to the right of the pattern).
    The orange bar shows the multiplicative speedup of our approach:
     $\EM / \GJ$.
    The blue bar shows the same, but \textit{excluding}
     the time spent building the indices needed for generic join:
     $\EM / (\GJ - \mathsf{idx})$.
    % Red ``$<$'' markers to the right of bars indicate that \EM timed out after 60 seconds.
    The text above each group of bars shows
     the pattern itself
     and
     the number of substitutions found on the largest \egraph;
     the patterns are sorted by this quantity.
  }\label{fig:speedup}
\end{figure}

\begin{table}
  \centering
  \caption{
    Summary statistics across patterns for each benchmark suite.
    The ``Idx'' column shows whether the time to build indices in \GJ
     is included (+) or not (--);
     the row color corresponds to the colors in \autoref{fig:speedup}.
    The ``Suite'' columns shows the benchmark suite,
     and the ``EG Size'' shows the number of \enodes
     in the \egraph used to benchmark.
    The ``\GJ'' and ``\EM'' columns show how many patterns
     that algorithm was fastest on in this configuration.
    % ``TO'' shows how many times \EM timed out
    %   (after 60 seconds); \GJ never timed out.
    ``Total'' shows the cumulative speedup across all patterns
      in that configuration.
    The remaining columns show statistics about the \EM/\GJ ratios
      for each pattern: harmonic mean, geometric mean, max, median, and min.
  }
  \rowcolors{2}{blue!20}{orange!20}
  \begin{tabular}{|llrrrrrrrrr|}
    \hline
    Idx  & Suite  & EG Size & \GJ & \EM & Total & HMean & GMean & Best  & Medn  & Worst   \\
    \hline
    +     & lambda & 4,142   & 15 & 3  & 1.69   & .84   & 1.71  & 13.62   & 1.60  & .12   \\
    --    & lambda & 4,142   & 18 & 0  & 2.58   & 2.99  & 4.23  & 39.17   & 3.68  & 1.10  \\
    +     & lambda & 57,454  & 16 & 2  & 2.60   & .95   & 2.66  & 136.54  & 2.65  & .12   \\
    --    & lambda & 57,454  & 18 & 0  & 2.87   & 3.33  & 9.11  & 406.70  & 4.05  & 1.03  \\
    +     & lambda & 109,493 & 15 & 3  & 1.66   & 1.75  & 3.11  & 148.96  & 2.03  & .65   \\
    --    & lambda & 109,493 & 18 & 0  & 1.70   & 3.32  & 7.46  & 291.18  & 4.10  & 1.05  \\
    +     & lambda & 213,345 & 15 & 3  & 2.20   & 1.55  & 3.40  & 304.33  & 1.72  & .43   \\
    --    & lambda & 213,345 & 18 & 0  & 2.21   & 2.96  & 8.23  & 501.12  & 5.04  & 1.04  \\
    \hline \rowcolor{white} \multicolumn{11}{c}{} \\
    \hline \rowcolor{white}
    Idx & Suite & EG Size & \GJ & \EM & Total  & HMean & GMean & Best         & Medn  & Worst \\
    \hline
    +   & math  & 8,205   & 30  & 2   & 5.49   & 0.64  & 4.61  & 66.54        & 2.79  & 0.03  \\
    --  & math  & 8,205   & 30  & 2   & 5.21   & 2.93  & 8.62  & 1,630.00     & 5.48  & 0.62  \\
    +   & math  & 53,286  & 29  & 3   & 311.23 & 2.61  & 13.50 & 50,030.29    & 3.62  & 0.74  \\
    --  & math  & 53,286  & 30  & 2   & 318.95 & 3.39  & 29.60 & 1,325,802.56 & 30.72 & 0.74  \\
    +   & math  & 132,080 & 29  & 3   & 96.55  & 2.66  & 15.18 & 61,488.73    & 4.02  & 0.60  \\
    --  & math  & 132,080 & 30  & 2   & 97.84  & 3.46  & 34.16 & 2,447,939.38 & 68.71 & 0.75  \\
    +   & math  & 217,396 & 30  & 2   & 119.82 & 2.83  & 18.34 & 101,023.37   & 3.91  & 0.72  \\
    --  & math  & 217,396 & 31  & 1   & 119.73 & 3.45  & 41.35 & 8,575,830.58 & 80.84 & 0.76  \\
    \hline
  \end{tabular}
  \vspace{1em}
  \label{tab:eval}
\end{table}

\subsection{Results}

\autoref{fig:speedup}
 show the results of our benchmarking experiments.
\GJ can be over 6 orders of magnitude faster
 than traditional e-matching on complex patterns.
Speedup tends to be greater when the output size is smaller,
 and when the pattern is larger and non-linear.
A large output indicates the \egraph\ is densely populated
 with terms matching the given pattern,
 therefore backtracking search wastes little time on unmatched terms,
 and using relational \ematching contributes little or no speedup.
Large and complex (non-linear) patterns require careful query planning
 to be processed efficiently.
For example, the pattern experiencing the largest speedup in~\autoref{fig:speedup}
 is 4 \enodes\ deep with 4 occurrences of the variable $f$.
Relational \ematching\ using generic join can devise a variable ordering
 to put smaller relations with fewer children on the outer loop,
 thereby pruning down a large search space early.
In contrast, backtracking search must traverse the \egraph\
 top-down.

In some cases index building time takes a significant proportion of the run time
 in relational \ematching, sometimes offsetting the gains.
Overall, relational \ematching\ remains competitive with the index building overhead.
In~\autoref{sec:incremental} we discuss potential remedies to alleviate this overhead.


\autoref{tab:eval}
 shows summary statistics across all patterns for each benchmark configuration.
Notably, \GJ is faster across patterns in
 every benchmarking configuration (the ``Total'' column).
 % the ``Total'' column shows that \GJ is faster overall than \EM
 % for every benchmarking configuration.
Much of the total benchmarking run time is dominated by
 simple, linear patterns (e.g. \texttt{(+ (+ a b) c)})
 that return many results.
Terms matching such patterns come to dominate
 the \egraph\ over time, due to the explosive expansion of
 associativity and commutativity.
As a result, the total speedup does not necessarily increase
 as the \egraph\ grows,
 whereas the best speedup as well as different average statistics
 steadily increase.

In summary, relational \ematching\ is almost always faster than backtracking search,
and is especially effective in speeding up complex patterns.

%%% Local Variables:
%%% mode: latex
%%% TeX-master: "main"
%%% End:
